% ============================================================================
\section{Example III: donor--acceptor electron transfer}
\label{sec:et}
% ============================================================================

\subsection{Diabatic states and a nuclear coordinate}

Consider an electron transferring from a donor to an acceptor,
$\Dm \longrightarrow \Am$. The two basis states are the \textbf{diabatic}
(charge-localized) states $\ket{D}\equiv\ket{\alpha}$ and $\ket{A}\equiv\ket{\beta}$.

What is new here is that their energies are not fixed numbers. Solvent molecules
reorganize around whichever site currently holds the charge, so the diabatic energies
depend on where the nuclei are. Introduce a single collective nuclear coordinate $q$ ---
the solvent or reaction coordinate --- and model each diabatic surface as a harmonic
well of the same force constant $k$, displaced by $\pm q_0$:
%
\begin{equation}
V_D(q) = \half k(q+q_0)^2,
\qquad
V_A(q) = \half k(q-q_0)^2 + \Delta G^{\circ}.
\label{eq:diabats}
\end{equation}
%
The donor state is happiest at $q=-q_0$, the acceptor state at $q=+q_0$, and
$\Delta G^\circ$ is the thermodynamic driving force. The electronic coupling
$V_{DA} = \braket{D|\Hop|A}$ between them is essentially independent of $q$.

\subsection{Identifying $\dvec$}

With $\Hop(q)$ built from \cref{eq:diabats},
%
\begin{equation}
d_z(q) = V_D(q)-V_A(q) = 2kq_0\,q - \Delta G^\circ,
\qquad
d_x = 2V_{DA},
\qquad d_y = 0.
\end{equation}
%
Defining the \textbf{reorganization energy} $\lambda \equiv 2kq_0^2$ --- the energy cost
of distorting the nuclei from the donor's optimal geometry to the acceptor's
\emph{without} moving the electron:

\begin{keyresult}[Electron transfer]
\vspace{-0.35cm}
\begin{equation}
d_z(q) = \frac{\lambda}{q_0}\,q - \Delta G^\circ,
\qquad
d_x = 2V_{DA}.
\label{eq:et-d}
\end{equation}
\end{keyresult}

\noindent
This is the structural heart of the example, and it is worth stating plainly:
\textbf{the $\szo$ coefficient is linear in a nuclear coordinate, and the $\sxo$
coefficient is a constant.} Compare \cref{eq:dimer-d}, where $d_z$ was a fixed
electronegativity difference. Here the molecule sweeps its own $d_z$ back and forth as
it vibrates, so the parameter scan of \cref{fig:avoided} is not a hypothetical --- it is
something the molecule physically does, many times per picosecond.

\subsection{Adiabatic surfaces and the Marcus barrier}

The eigenvalues of $\Hop(q)$ are the \textbf{adiabatic} potential energy surfaces --- the
Born--Oppenheimer curves the nuclei actually move on:
%
\begin{equation}
E_\pm(q) = \bar{V}(q) \pm \half\sqrt{d_z(q)^2 + 4V_{DA}^2},
\qquad \bar{V}(q) = \half\!\left[V_D(q)+V_A(q)\right].
\label{eq:adiabats}
\end{equation}
%
Three consequences follow at once from \cref{sec:geometry}:

\begin{itemize}[leftmargin=*,itemsep=2pt]
\item the diabats cross where $d_z(q^\ddagger)=0$, that is at
      $q^\ddagger = \Delta G^\circ q_0/\lambda$;
\item the adiabatic gap there is exactly $|d_x| = 2V_{DA}$;
\item the activation energy, measured on the \emph{diabatic} surface from the donor
      minimum up to the crossing, is
      \begin{equation}
      \Delta G^{\ddagger} = V_D(q^\ddagger) - V_D(-q_0)
          = \frac{(\lambda+\Delta G^\circ)^2}{4\lambda},
      \label{eq:marcus-barrier}
      \end{equation}
      the \textbf{Marcus} activation free energy \cite{Marcus1956,MarcusNobel}
      (algebra in \cref{app:marcus}).
\end{itemize}

\Cref{eq:marcus-barrier} is a parabola in the driving force, minimized --- and equal to
zero --- at $\Delta G^\circ = -\lambda$. Since the nonadiabatic rate goes as
$k_{ET}\propto |V_{DA}|^2\exp(-\Delta G^\ddagger/k_BT)$, making a reaction \emph{more}
exergonic speeds it up only until $-\Delta G^\circ = \lambda$, and slows it down
thereafter. This is the celebrated \textbf{inverted region}, whose experimental
confirmation was central to the 1992 Nobel Prize in Chemistry. \Cref{fig:et,fig:marcus}
show the surfaces and the rate.

% ---------------------------------------------------------------------------
\begin{figure}[t]
\centering
\begin{tikzpicture}
\begin{axis}[
  widenote, height=7.4cm,
  xlabel={nuclear (solvent) coordinate $q$},
  ylabel={energy},
  xmin=-2.25, xmax=2.25, ymin=-0.95, ymax=1.75,
  legend style={at={(0.5,-0.22)},anchor=north,legend columns=2,draw=none,fill=none,
                font=\footnotesize},
  legend cell align=left,
]
% adiabats: soft thick bands underneath
\addplot[cLower, line width=3.6pt, opacity=0.35, domain=-2.25:2.25, samples=400]
  {0.5*(0.5*(x+1)^2 + 0.5*(x-1)^2 - 0.5) - 0.5*sqrt((2*x+0.5)^2 + 0.09)};
\addlegendentry{adiabat $E_-$}
\addplot[cUpper, line width=3.6pt, opacity=0.35, domain=-2.25:2.25, samples=400]
  {0.5*(0.5*(x+1)^2 + 0.5*(x-1)^2 - 0.5) + 0.5*sqrt((2*x+0.5)^2 + 0.09)};
\addlegendentry{adiabat $E_+$}
% diabats on top, crisp
\addplot[cMuted, dashed, thick, domain=-2.25:2.25, samples=200]{0.5*(x+1)^2};
\addlegendentry{diabat $V_D$ (charge on donor)}
\addplot[cMuted, densely dotted, thick, domain=-2.25:2.25, samples=200]{0.5*(x-1)^2-0.5};
\addlegendentry{diabat $V_A$ (charge on acceptor)}
% avoided crossing marker at q* = -0.25
\draw[<->,cAccent,thick] (axis cs:-0.25,0.13125) -- (axis cs:-0.25,0.43125);
\node[cAccent,font=\footnotesize,anchor=west] at (axis cs:-0.17,0.60){$2V_{DA}$};
% Marcus barrier
\draw[<->,cWarm,thick] (axis cs:-1.28,0.0) -- (axis cs:-1.28,0.2812);
\draw[cWarm,densely dash dot,thin] (axis cs:-1.36,0.2812) -- (axis cs:-0.25,0.2812);
\draw[cWarm,densely dash dot,thin] (axis cs:-1.36,0.0) -- (axis cs:-1.0,0.0);
\node[cWarm,font=\footnotesize,anchor=east,fill=white,inner sep=1pt]
  at (axis cs:-1.34,0.145){$\Delta G^{\ddagger}$};
% minima
\addplot[only marks, mark=*, mark size=1.6pt, black] coordinates {(-1,0) (1,-0.5)};
\node[font=\footnotesize,anchor=north] at (axis cs:-1,-0.07){$\Dm$};
\node[font=\footnotesize,anchor=north] at (axis cs:1,-0.57){$\Am$};
% driving force
\draw[<->,cPurple,thin] (axis cs:1.62,0.0) -- (axis cs:1.62,-0.5);
\node[cPurple,font=\footnotesize,anchor=west] at (axis cs:1.66,-0.25){$\Delta G^{\circ}$};
\draw[cPurple,densely dotted,thin] (axis cs:-1,0) -- (axis cs:1.62,0);
\end{axis}
\end{tikzpicture}
\caption{Diabatic and adiabatic surfaces for
$k=1$, $q_0=1$ (so $\lambda=2$), $\Delta G^\circ = -0.5$, $V_{DA}=0.15$. Away from the
crossing the adiabats and diabats coincide --- the electron stays put. Near
$q^\ddagger = \Delta G^\circ q_0/\lambda = -0.25$ the coupling splits them by $2V_{DA}$
and the character of each adiabatic surface switches over.}
\label{fig:et}
\end{figure}
% ---------------------------------------------------------------------------

% ---------------------------------------------------------------------------
\begin{figure}[t]
\centering
\begin{tikzpicture}
\begin{groupplot}[
  group style={group size=2 by 1, horizontal sep=1.5cm},
  notestyle, height=5.8cm,
]
\nextgroupplot[
  xlabel={$\Delta G^\circ/\lambda$}, ylabel={$\Delta G^{\ddagger}/\lambda$},
  title={Activation barrier},
  xmin=-3, xmax=0.5, ymin=-0.03, ymax=1.05,
]
\addplot[cWarm, very thick, domain=-3:0.5, samples=200]{(1+x)^2/4};
\draw[cMuted,densely dotted,thin] (axis cs:-1,-0.03) -- (axis cs:-1,1.05);
\node[cMuted,font=\footnotesize,align=center] at (axis cs:-2.1,0.62){inverted\\region};
\node[cMuted,font=\footnotesize,align=center] at (axis cs:-0.35,0.62){normal\\region};

\nextgroupplot[
  xlabel={$\Delta G^\circ/\lambda$}, ylabel={rate $\propto e^{-\Delta G^{\ddagger}/k_BT}$},
  title={Marcus rate ($k_BT=\lambda/20$)},
  xmin=-3, xmax=0.5, ymode=log, ymin=1e-9, ymax=4,
  ytick={1e-8,1e-6,1e-4,1e-2,1},
]
\addplot[cPurple, very thick, domain=-3:0.5, samples=300]{exp(-5*(1+x)^2)};
\draw[cMuted,densely dotted,thin] (axis cs:-1,1e-9) -- (axis cs:-1,4);
\end{groupplot}
\end{tikzpicture}
\caption{The Marcus inverted region. The barrier \eqref{eq:marcus-barrier} is a parabola
in the driving force, so the rate is \emph{not} monotonic: past
$-\Delta G^\circ=\lambda$ a more exergonic reaction is a slower one.}
\label{fig:marcus}
\end{figure}
% ---------------------------------------------------------------------------

\subsection{Landau--Zener: sweeping through the crossing}

Everything above is static --- surfaces evaluated at fixed $q$. Now let the nuclei
actually move through the crossing. If the coordinate sweeps linearly in time, $q = vt$,
then by \cref{eq:et-d}
%
\begin{equation}
d_z(t) = \alpha t, \qquad \alpha = \frac{\lambda v}{q_0}
  = \left|\frac{\dd}{\dd t}\big(V_D-V_A\big)\right| ,
\end{equation}
%
the rate at which the diabatic gap is swept. Starting far to one side in a single
diabatic state, two outcomes compete:

\begin{itemize}[leftmargin=*,itemsep=2pt]
\item \textbf{adiabatic passage} --- the system follows the lower adiabatic surface and
therefore \emph{changes} diabatic character: the electron transfers;
\item \textbf{diabatic passage} --- the system blasts straight through the avoided
crossing, staying on the same diabat: no transfer.
\end{itemize}

\noindent
Landau and Zener solved this exactly in 1932 \cite{Zener1932,Landau1932}:

\begin{keyresult}[Landau--Zener]
\vspace{-0.35cm}
\begin{equation}
P_{\rm diabatic} = \exp\!\left(-\frac{2\pi V_{DA}^2}{\hbar\,\alpha}\right).
\label{eq:lz}
\end{equation}
\end{keyresult}

\noindent
Slow sweeps or strong coupling give adiabatic passage ($P\to0$); fast sweeps or weak
coupling give diabatic passage ($P\to1$). Note the combination in the exponent:
$(\text{coupling})^2$ divided by (sweep rate). Expanding for small $V_{DA}$,
%
\begin{equation}
1 - P_{\rm diabatic} \approx \frac{2\pi V_{DA}^2}{\hbar\alpha},
\end{equation}
%
which is the same $|V_{DA}|^2$ proportionality that appears in Fermi's golden rule and
in the nonadiabatic Marcus rate --- the two results are the same physics seen from two
directions.

% ---------------------------------------------------------------------------
\begin{figure}[t]
\centering
\begin{tikzpicture}
\begin{axis}[
  notestyle, width=0.62\textwidth, height=6.2cm,
  xlabel={diabatic coupling $V_{DA}$},
  ylabel={$P_{\rm diabatic}$},
  xmin=0, xmax=0.53, ymin=0.13, ymax=1.05,
  legend pos=north east, legend cell align=left,
]
\addplot[black, very thick, domain=0:0.53, samples=200]{exp(-2*pi*x^2)};
\addlegendentry{Zener formula, Eq.~(\ref{eq:lz})}
\addplot[only marks, mark=*, mark size=2.1pt, cUpper]
  table[x index=0, y index=1] {figures/landau_zener_numeric.dat};
\addlegendentry{numerical propagation}
\end{axis}
\end{tikzpicture}
\caption{\Cref{eq:lz} against direct numerical propagation through the crossing
($\hbar=1$, sweep rate $\alpha=1$). The numerical points come from the notebook, which
slices the sweep into 60000 steps and evaluates $\Hop$ at the midpoint of each. The
residual disagreement, at most $0.004$, is not step error: Zener's result is exact only
for a sweep running from $t=-\infty$ to $+\infty$, and any truncated sweep leaves a
little interference behind.}
\label{fig:lz}
\end{figure}
% ---------------------------------------------------------------------------

\begin{lookahead}[Where this is going: the Holstein model]
Collecting \cref{eq:et-d},
\begin{equation}
\Hop(q) = \bar{V}(q)\,\Iop
 + \half\!\left[\frac{\lambda}{q_0}q - \Delta G^\circ\right]\szo + V_{DA}\,\sxo .
\end{equation}
We treated $q$ as a \emph{classical parameter} --- a number to dial, or to sweep linearly
in time. But $q$ is a nuclear coordinate, and nuclei are quantum. Promote it to a
harmonic oscillator, $\hat{q}\propto(\hat{a}+\hat{a}^\dagger)$, and
\begin{equation}
\Hop_{\rm Holstein} = \frac{\varepsilon}{2}\szo + V_{DA}\,\sxo
  + \hbar\omega\,\hat{a}^\dagger\hat{a}
  + \hbar\omega g\,(\hat{a}+\hat{a}^\dagger)\,\szo ,
\label{eq:holstein}
\end{equation}
the \textbf{Holstein}, or spin--boson, Hamiltonian \cite{Holstein1959,Leggett1987}.

Set \cref{eq:holstein} beside \cref{eq:jc} and notice how different they are despite
sharing a two-dimensional electronic space and one boson:
\begin{center}
\small
\begin{tabular}{@{}lll@{}}
\toprule
Model & Coupling operator & The boson\ldots \\
\midrule
Jaynes--Cummings & $\hat{a}\hat\sigma_+ + \hat{a}^\dagger\hat\sigma_-$
  & \textbf{flips} the two-level system \\
Holstein & $(\hat{a}+\hat{a}^\dagger)\,\szo$
  & \textbf{shifts} its energies \\
\bottomrule
\end{tabular}
\end{center}
One gives polaritons and vacuum Rabi splitting; the other gives vibronic progressions,
Franck--Condon factors, and polaron formation. You have now built the single-spin
ancestor of both.
\end{lookahead}
